\section{Introduction}\label{sec:tsps-introduction}

Structure-preserving
signatures~\cite{abe2010structure,Abe2011,Abe2014,Kiltz2015,Groth2015,Ghadafi2016}
are digital signatures in which the messages, public keys, and signatures are
all group elements, and verification consists solely of checking pairing
product equations. Coupled with non-interactive zero-knowledge proofs,
structure-preserving signatures have proved useful across a range of
privacy-preserving applications: blind signatures~\cite{Fuchsbauer2010},
delegatable anonymous credentials~\cite{Camenisch2017,Crites2019,Mir2023}, and
group signatures~\cite{Groth2006,Groth2007,Libert2015}, to name a few.

Structure-preserving signatures also find use in blockchain-based
privacy-preserving payment systems. In~\cite{Androulaki2020,Tomescu2022,Kiayias2022},
currency is represented by a token encoding the triple $\langle
\text{owner}, \ \text{value}, \ \text{serial\_number} \rangle$. For privacy,
the shared ledger records a hiding commitment to the token, generally a
Pedersen commitment~\cite{Pedersen1991}, rather than the token itself. To
transfer ownership privately, the owner interacts with a certifier to obtain a
blind Pointcheval-Sanders
signature~\cite{Pointcheval2016,Pointcheval2018,Sonnino2019} on the
commitment's opening, which is later used to prove that the token being spent
is recorded on the ledger, and, via its serial number, that it has not been
spent before. Structure-preserving signatures remove the need for this
interaction: since signing acts directly on the group elements that make up
the commitment, the certifier can sign tokens as they appear on the ledger and
simply publish the resulting signatures.

A similar idea applies where the ledger stores commitments to user accounts,
pairs $\langle \text{owner}, \ \text{holdings} \rangle$, rather than committed
tokens. Such accounts form the basis for periodic reporting on user holdings,
for instance proving that holdings do not exceed some threshold when filing
taxes. Since tax authorities typically lack direct access to the shared
ledger, they instead rely on signatures produced by the ledger's
administrators. Structure-preserving signatures let the administrators sign
the committed accounts together with a timestamp and publish the result
directly, without involving the user. The user can then use the published
signature to prove properties of their holdings without disclosing anything
beyond what is required.

In both applications, relying on a centralised certifier or administrator sits
uneasily with the decentralised nature of blockchain systems, so
structure-preserving signatures are only fit for this purpose once they admit
an efficient threshold counterpart. We give two threshold
structure-preserving signature schemes, building on the signatures
of~\cite{Abe2011} and~\cite{Kiltz2015}, together with the secure multiparty
computation (MPC) techniques
of~\cite{damgaard2012multiparty,keller2016mascot,Goyal2021}. We select
\cite{Abe2011} for its optimal signature size, and \cite{Kiltz2015} for its
reliance on standard assumptions alone. Both underlying signatures let the
threshold signers perform the bulk of their computation offline, in a
pre-processing phase before the message to be signed is known, which suits
high-throughput applications, such as payments, that experience interspersed
periods of low and high traffic. Since our threshold signatures require
distributed computation over group elements rather than only field elements,
we build on the techniques of~\cite{Ghosh2023}, which extend MPC protocols
over field elements to protocols over group elements. Our schemes inherit the
security properties of the underlying structure-preserving signatures, and
the cost of the online phase of threshold signing is comparable to that of
ordinary signing.

\subsection{Related work}\label{sec:tsps-related}

\paragraph{Structure-preserving signatures.} Structure-preserving signatures
have messages, signatures, and verification keys that are all group elements.
They were first introduced in~\cite{abe2010structure} and have since been
studied extensively, with matching impossibility results known for signature
size in the generic group model (GGM), from $q$-type assumptions, and from
non-$q$-type assumptions~\cite{Abe2011,abe2011separating}.

The scheme of~\cite{abe2010structure} signs messages of arbitrary length with
signatures of seven group elements, and is secure in the standard model under
a non-interactive assumption. \cite{Abe2011} then established a lower bound of
three group elements for signature size in type III pairing groups, those
without an efficiently computable isomorphism between $\GroupOne$ and
$\GroupTwo$, and gave a matching construction secure in the GGM. These
signatures achieve strong unforgeability, and remain three group elements
regardless of message length. \cite{libert2013linearly} gives a linearly
homomorphic structure-preserving signature, also of just three group elements,
enabling further applications.

Fully structure-preserving signatures additionally have secret keys that are
group elements~\cite{abe2019efficient}. Since structure-preserving signatures
integrate well with Groth-Sahai proofs of knowledge~\cite{groth2008efficient},
a group-element secret key lets a user also prove statements about their own
key; the cost is that signature length grows with the square root of the
message length.

\paragraph{Threshold structure-preserving signatures.} Threshold
structure-preserving signatures have only recently been considered.
\cite{crites2023threshold} gives a non-interactive construction, but one
restricted to signing indexed Diffie-Hellman tuples of the form $(\genOne^m,
\genTwo^m)$. This restricts the scheme to messages whose discrete logarithm
relative to $\genOne$ or $\genTwo$ is known, and rules out signing an existing
public key or commitment directly: either the threshold signers would learn
someone else's secret key, or the element simply could not be signed under the
scheme. \cite{Mitrokotsa2024} instead gives a non-interactive threshold
structure-preserving signature capable of signing arbitrary group elements,
secure under the SXDH assumption. It builds on the structure-preserving
signature of Kiltz, Pan, and Wee~\cite{Kiltz2015}, and its signatures consist
of seven group elements.

We give two threshold structure-preserving signature schemes for arbitrary
group elements, built respectively from~\cite{Abe2011} and~\cite{Kiltz2015},
that minimise the computation and communication required of parties in the
online phase. Table~\ref{table:comp} in Section~\ref{sec:eval} compares the
resulting schemes with~\cite{crites2023threshold} and~\cite{Mitrokotsa2024} in
terms of assumptions, key and signature size, and verification cost.

\paragraph{Our contribution.} The contributions of this chapter are threefold.

\begin{itemize}

  \item We design two threshold structure-preserving signature schemes in the
    hybrid model: the constructions assume the existence of ideal
    functionalities that model the computations of the threshold signers.

  \item We prove the security of the threshold signatures in the hybrid
    model, and show that they offer the same security guarantees as their
    non-threshold counterparts.

  \item We benchmark the threshold signatures using two MPC frameworks that
    realise these ideal functionalities: ATLAS, which assumes an honest
    majority~\cite{Goyal2021}, and MASCOT, which assumes a dishonest
    majority~\cite{keller2016mascot}. The former suits settings where
    robustness and high availability are desirable, the latter settings
    where resilience to key compromise takes precedence.

\end{itemize}

We design our schemes to require no interaction between parties during the
online phase, and, to the best of our knowledge, give the first multiparty
secret sharing of group elements, building on prior work extending secret
sharing to group elements~\cite{Ghosh2023}. We also implement previously
suggested threshold structure-preserving signature schemes, so as to compare
concrete efficiency. Our scheme built from~\cite{Abe2011} yields optimally
small threshold signatures, consisting of just three group elements.

\paragraph{Layout.} Section~\ref{sec:formalization} formalises the security
properties of threshold structure-preserving signatures.
Section~\ref{sec:SPS} describes the structure-preserving signatures of Abe,
Groth, Haralambiev, and Ohkubo~\cite{Abe2011} and of Kiltz, Pan, and
Wee~\cite{Kiltz2015}. Section~\ref{sec:MPC} introduces the ideal
functionalities for MPC that our constructions build on.
Section~\ref{sec:schemes} shows how these functionalities threshold sign with
both signature schemes, proven secure in Section~\ref{sec:tsps-security}.
Section~\ref{sec:eval} evaluates our approach, realising the MPC
functionalities concretely via ATLAS and MASCOT, and
Section~\ref{sec:tsps-summary} concludes.
